\section{Limitations and open points}\label{sec:open}

The following points are either not verified or rest on declared choices; they are listed here
rather than stated as results.

\begin{itemize}
\item \emph{Scope of the model.} All simulations are two-dimensional, with a fixed $N$
(multinomial statistics), a background equal in all exposures with the SBR held fixed at every
position (Sec.~\ref{sec:model}), no dipole orientation effects,
no aberrations beyond the zero depth and no drift. The camera reference is ideal
(Sec.~\ref{sec:camera}); real cameras are worse.
\item \emph{rmse normalization.} We use the per-axis convention of Eq.~\eqref{eq:crb} for $\sigma$
and rmse. Whether it coincides with the error definition of Masullo \textit{et al.}
\cite{Masullo2022} was not checked, because the main text of that work is not in our notes.
\item \emph{$\sbr=10$ for the MLE efficiency.} The efficiency at the centre is defined with
$\sbr=10$; this was chosen by the study coordinator on behalf of the author because without
background the model is not regular (Sec.~\ref{sec:estimators}). It is a convention, pending
explicit confirmation by the author.
\item \emph{Vectorial ring diameter versus PyFocus.} Our peak-to-peak diameters of the
vectorial donut as a function of the filling factor ($n=1.5$) agree to within a few percent with
the values we transcribed from Fig.~10 of Ref.~\cite{Caprile2022} at filling factors $0.5$, $2$
and uniform illumination, but differ clearly at filling factor $1$. We do not claim agreement
there; an independent comparison is pending.
\item \emph{Heavy tail at low photon numbers.} In iterative MINFLUX at $N_{\mathrm{tot}}=250$
the error distribution is heavy-tailed and the bootstrap SE of $\sigma$ is itself unstable; the
full-range slope inherits this.
\item \emph{Range dependence of the $L$ exponent.} The exponent
\pnum{crb_exponent_L}\src{crb_exponent_L} applies only to $L\ll\fwhm$; it increases when the fit
range extends to larger $L$.
\item \emph{Three-dimensional donuts and axial misalignment} are outside the scope of this study.
\item \emph{Instrument-level effects} studied with SimuFLUX \cite{MarinRies2026}, i.e.\
misalignment of the vortex phase mask, aberrations, fluorophore blinking and bleaching,
neighbouring emitters, emitter motion, unequal dwell times of the exposures and the validity
filters of commercial instruments, are not modelled here. In particular all our exposures have
the same dwell time.
\end{itemize}
